% =========================================================
% MODULO DE PERSISTENCIA (BASE)
% =========================================================

\subsection{Modelo de Persistencia Base}
Se define el esquema de la base de datos relacional utilizando PostgreSQL. El diseño sigue el paradigma de mapeo objeto-relacional (ORM) a trav\'es de JPA/Hibernate, garantizando la consistencia y escalabilidad del modelo de datos.

\subsubsection{Diccionario de Datos: Tabla Usuarios}
La tabla \texttt{usuarios} constituye la entidad central para la autenticaci\'on y gesti\'on de perfiles dentro del sistema.

\begin{table}[H]
\centering
\small
\begin{tabularx}{\textwidth}{|l|l|l|X|}
\hline
\rowcolor{azulESCOM} 
\color{white}\textbf{Campo} & \color{white}\textbf{Tipo de Dato} & \color{white}\textbf{Restricciones} & \color{white}\textbf{Descripci\'on} \\ \hline
id & BIGINT & PK, IDENTITY & Identificador \'unico interno del registro. \\ \hline
correo & VARCHAR(255) & UNIQUE, NOT NULL & Correo institucional del usuario. \\ \hline
password & VARCHAR(255) & NOT NULL & Hash criptogr\'afico de la contrase\~na de acceso. \\ \hline
boleta & VARCHAR(50) & UNIQUE, NOT NULL & Identificador institucional (boleta del alumno o empleado). \\ \hline
nombre & VARCHAR(255) & - & Nombre(s) del titular del perfil. \\ \hline
primer\_apellido & VARCHAR(255) & - & Apellido paterno. \\ \hline
segundo\_apellido & VARCHAR(255) & - & Apellido materno (campo opcional). \\ \hline
reset\_token & VARCHAR(255) & UNIQUE & Token empleado para la recuperaci\'on de la cuenta. \\ \hline
token\_expiration & TIMESTAMP & - & Fecha y hora de expiraci\'on del token de seguridad. \\ \hline
\end{tabularx}
\caption{Diccionario de datos de la entidad \texttt{usuarios}.}
\label{tab:dict_usuarios}
\end{table}

\subsubsection{Diccionario de Datos: Tabla Usuario Folios}
Tabla auxiliar que gestiona la relaci\'on din\'amica de folios vinculados a cada usuario.

\begin{table}[H]
\centering
\small
\begin{tabularx}{\textwidth}{|l|l|l|X|}
\hline
\rowcolor{azulESCOM} 
\color{white}\textbf{Campo} & \color{white}\textbf{Tipo de Dato} & \color{white}\textbf{Restricciones} & \color{white}\textbf{Descripci\'on} \\ \hline
usuario\_id & BIGINT & FK, NOT NULL & Clave for\'anea que referencia a la tabla \texttt{usuarios}. \\ \hline
folio & VARCHAR(255) & PK & C\'odigo \'unico del folio registrado. \\ \hline
boleta & VARCHAR(50) & - & Boleta espec\'ificamente asociada a dicho folio. \\ \hline
\end{tabularx}
\caption{Diccionario de datos de la tabla auxiliar \texttt{usuario\_folios}.}
\label{tab:dict_usuario_folios}
\end{table}

\subsubsection{Definici\'on de la Entidad Usuario (JPA)}
\begin{lstlisting}[language=Java, caption={Entidad Java para la gesti\'on de usuarios.}, label={lst:java_usuario}]
@Entity
@Table(name = "usuarios")
@Data
@NoArgsConstructor
@AllArgsConstructor
public class UsuarioEntity {
    @Id
    @GeneratedValue(strategy = GenerationType.IDENTITY)
    private Long id;

    @Column(unique = true, nullable = false)
    private String correo;
    
    @Column(nullable = false)
    private String password;
    
    @Column(unique = true, nullable = false, length = 50)
    private String boleta;
    
    private String nombre;
    private String primerApellido;
    private String segundoApellido;
    
    @Column(unique = true)
    private String resetToken;
    
    private LocalDateTime tokenExpiration;
}
\end{lstlisting}

\subsubsection{Definici\'on del Esquema SQL (M\'odulo Base)}
El siguiente script define la estructura f\'isica en PostgreSQL para el m\'odulo de usuarios.

\begin{lstlisting}[language=SQL, caption={Script de creaci\'on para el m\'odulo de usuarios.}, label={lst:sql_usuarios}]
CREATE TABLE usuarios (
    id BIGSERIAL PRIMARY KEY,
    correo VARCHAR(255) UNIQUE NOT NULL,
    password VARCHAR(255) NOT NULL,
    boleta VARCHAR(50) UNIQUE NOT NULL,
    nombre VARCHAR(255),
    primer_apellido VARCHAR(255),
    segundo_apellido VARCHAR(255),
    reset_token VARCHAR(255) UNIQUE,
    token_expiration TIMESTAMP
);

CREATE TABLE usuario_folios (
    usuario_id BIGINT NOT NULL,
    folio VARCHAR(255) NOT NULL,
    boleta VARCHAR(50),
    PRIMARY KEY (usuario_id, folio),
    FOREIGN KEY (usuario_id) REFERENCES usuarios(id) ON DELETE CASCADE
);


CREATE INDEX idx_usuario_folios_usuario_id ON usuario_folios(usuario_id);
\end{lstlisting}

% =========================================================
% MODULO DE NOTIFICACIONES
% =========================================================

\subsection{M\'odulo de Notificaciones}
Este m\'odulo extiende el modelo de persistencia base para gestionar el almacenamiento de alertas y comunicaciones dirigidas a los usuarios. El diseño permite mantener un hist\'orico completo de mensajes vinculados a cada perfil, facilitando el seguimiento de actividades y avisos del sistema.

\subsubsection{Diccionario de Datos: Tabla Notificaciones}
La tabla \texttt{notificaciones} almacena los mensajes generados por el sistema, incorporando mecanismos para el rastreo de lectura y la vinculaci\'on con el usuario receptor.

\begin{table}[H]
\centering
\small
\begin{tabularx}{\textwidth}{|l|l|l|X|}
\hline
\rowcolor{azulESCOM} 
\color{white}\textbf{Campo} & \color{white}\textbf{Tipo de Dato} & \color{white}\textbf{Restricciones} & \color{white}\textbf{Descripci\'on} \\ \hline
id & BIGINT & PK, IDENTITY & Identificador \'unico de la notificaci\'on. \\ \hline
titulo & VARCHAR(255) & NOT NULL & Encabezado o asunto del mensaje. \\ \hline
mensaje & TEXT & NOT NULL & Contenido detallado de la notificaci\'on. \\ \hline
boleta & VARCHAR(50) & - & Boleta asociada al evento (redundancia para optimizaci\'on de consultas). \\ \hline
fecha & TIMESTAMP & NOT NULL & Fecha y hora de creaci\'on del registro (por defecto la actual). \\ \hline
leida & BOOLEAN & DEFAULT FALSE & Estado que indica si la notificaci\'on ha sido visualizada. \\ \hline
usuario\_id & BIGINT & FK, NOT NULL & Referencia al titular en la tabla \texttt{usuarios}. \\ \hline
\end{tabularx}
\caption{Diccionario de datos de la entidad \texttt{notificaciones}.}
\label{tab:dict_notificaciones}
\end{table}

\subsubsection{Definici\'on de la Entidad (JPA)}
A continuaci\'on se presenta la implementaci\'on de la clase de entidad en Java. Se emplea la librer\'ia \textbf{Lombok} mediante la anotaci\'on \texttt{@Data} para la generaci\'on autom\'atica de m\'etodos de acceso (getters/setters), constructores y otros utilitarios.

\begin{lstlisting}[language=Java, caption={Entidad Java para la gesti\'on de notificaciones.}, label={lst:java_notificacion}]
@Entity
@Table(name = "notificaciones")
@Data
@NoArgsConstructor
@AllArgsConstructor
public class NotificacionEntity {
    @Id
    @GeneratedValue(strategy = GenerationType.IDENTITY)
    private Long id;

    @Column(nullable = false)
    private String titulo;
    
    @Column(nullable = false, columnDefinition = "TEXT")
    private String mensaje;
    
    private String boleta;
    
    @Column(nullable = false)
    private LocalDateTime fecha;
    
    @Column(columnDefinition = "BOOLEAN DEFAULT FALSE")
    private boolean leida = false;

    @ManyToOne
    @JoinColumn(name = "usuario_id", nullable = false)
    private UsuarioEntity usuario;
}
\end{lstlisting}

\subsubsection{Definici\'on del Esquema SQL (M\'odulo Notificaciones)}
Script definitivo para la creaci\'on de la estructura f\'isica en PostgreSQL, incluyendo la restricci\'on de integridad referencial con borrado en cascada.

\begin{lstlisting}[language=SQL, caption={Script de creaci\'on para el m\'odulo de notificaciones.}, label={lst:sql_notificaciones}]
CREATE TABLE notificaciones (
    id BIGSERIAL PRIMARY KEY,
    titulo VARCHAR(255) NOT NULL,
    mensaje TEXT NOT NULL,
    boleta VARCHAR(50),
    fecha TIMESTAMP NOT NULL DEFAULT CURRENT_TIMESTAMP,
    leida BOOLEAN DEFAULT FALSE,
    usuario_id BIGINT NOT NULL,
    CONSTRAINT fk_usuario_notificacion 
        FOREIGN KEY (usuario_id) 
        REFERENCES usuarios(id) 
        ON DELETE CASCADE
);


CREATE INDEX idx_notificaciones_usuario_fecha ON notificaciones(usuario_id, fecha DESC);
\end{lstlisting}

% =========================================================
% MODULO DE FOLIOS (ENTIDAD MAESTRA)
% =========================================================

\section{M\'odulo de Gesti\'on de Folios}
La entidad \texttt{FolioEntity} constituye el n\'ucleo de informaci\'on de la plataforma. En ella se consolidan los datos del incidente, la informaci\'on de identificaci\'on del quejoso y, en cumplimiento con la normativa de protecci\'on a menores, los datos de los tutores legales cuando el caso lo requiere.

\subsection{Diccionario de Datos: Tabla Folios}
Debido a la naturaleza integral del formulario de denuncia, la tabla \texttt{folios} se divide en tres segmentos informativos: datos del expediente, datos personales y datos de contacto de terceros.

\begin{table}[H]
\centering
\small
\begin{tabularx}{\textwidth}{|l|l|l|X|}
\hline
\rowcolor{azulESCOM} 
\color{white}\textbf{Campo} & \color{white}\textbf{Tipo de Dato} & \color{white}\textbf{Restricciones} & \color{white}\textbf{Descripci\'on} \\ \hline
\multicolumn{4}{|l|}{\cellcolor[HTML]{EFEFEF}\textbf{Datos del Expediente}} \\ \hline
id & BIGINT & PK, IDENTITY & Identificador interno de la base de datos. \\ \hline
codigo\_folio & VARCHAR(255) & UNIQUE, NOT NULL & Folio p\'ublico generado para el seguimiento. \\ \hline
fecha\_creacion & TIMESTAMP & DEFAULT NOW() & Fecha y hora de registro en el sistema. \\ \hline
asunto & VARCHAR(255) & - & T\'itulo descriptivo de la queja. \\ \hline
fecha\_hechos & VARCHAR(255) & - & Fecha aproximada en que ocurrieron los hechos. \\ \hline
descripcion & TEXT & - & Descripci\'on detallada de los hechos denunciados. \\ \hline
unidad & VARCHAR(255) & - & Unidad acad\'emica o administrativa destino. \\ \hline
\multicolumn{4}{|l|}{\cellcolor[HTML]{EFEFEF}\textbf{Datos del Quejoso}} \\ \hline
boleta & VARCHAR(50) & - & N\'umero de boleta o empleado del titular. \\ \hline
nombre & VARCHAR(255) & - & Nombre(s) del denunciante. \\ \hline
primer\_apellido & VARCHAR(255) & - & Apellido paterno del denunciante. \\ \hline
segundo\_apellido & VARCHAR(255) & - & Apellido materno del denunciante. \\ \hline
correo & VARCHAR(255) & - & Correo electr\'onico de contacto. \\ \hline
\multicolumn{4}{|l|}{\cellcolor[HTML]{EFEFEF}\textbf{Representaci\'on Legal (Menores de edad)}} \\ \hline
nombre\_tutor & VARCHAR(255) & - & Nombre completo del tutor legal. \\ \hline
parentesco & VARCHAR(100) & - & Relaci\'on jur\'idica con el menor. \\ \hline
correo\_tutor & VARCHAR(255) & - & Correo de contacto del representante. \\ \hline
telefono\_tutor & VARCHAR(20) & - & Tel\'efono de contacto del tutor. \\ \hline
\end{tabularx}
\caption{Diccionario de datos de la entidad Folios (segmentado).}
\label{tab:dict_folios}
\end{table}

\subsection{Definici\'on de la Entidad (JPA)}
La clase \texttt{FolioEntity} gestiona la composici\'on de evidencias mediante una relaci\'on de uno a muchos, permitiendo que un expediente contenga m\'ultiples archivos probatorios con persistencia en cascada.

\begin{lstlisting}[language=Java, caption={Implementaci\'on de la entidad maestra de folios.}, label={lst:java_folio}]
@Entity
@Table(name = "folios")
@Data
@NoArgsConstructor
@AllArgsConstructor
public class FolioEntity {
    @Id
    @GeneratedValue(strategy = GenerationType.IDENTITY)
    private Long id;

    @Column(unique = true, nullable = false)
    private String codigoFolio;
    
    private LocalDateTime fechaCreacion;
    private String asunto;
    private String fechaHechos;
    private String descripcion;
    private String unidad;
    
    // Datos del quejoso
    private String nombre;
    private String primerApellido;
    private String segundoApellido;
    private String correo;
    private String boleta;
    
    // Datos del tutor
    private String nombreTutor;
    private String parentesco;
    private String correoTutor;
    private String telefonoTutor;
    
    @OneToMany(mappedBy = "folio", cascade = CascadeType.ALL, orphanRemoval = true)
    private List<EvidenciaEntity> evidencias = new ArrayList<>();
    
    
    public void addEvidencia(EvidenciaEntity evidencia) {
        evidencias.add(evidencia);
        evidencia.setFolio(this);
    }
    
    public void removeEvidencia(EvidenciaEntity evidencia) {
        evidencias.remove(evidencia);
        evidencia.setFolio(null);
    }
}
\end{lstlisting}

\subsection{Definici\'on del Esquema SQL}
El script de creaci\'on integra todas las dimensiones de la queja, asegurando la unicidad del c\'odigo de folio para evitar colisiones en el seguimiento p\'ublico.

\begin{lstlisting}[language=SQL, caption={Script DDL para la tabla maestra de folios.}, label={lst:sql_folios}]
CREATE TABLE folios (
    id BIGSERIAL PRIMARY KEY,
    codigo_folio VARCHAR(255) UNIQUE NOT NULL,
    fecha_creacion TIMESTAMP DEFAULT CURRENT_TIMESTAMP,
    asunto VARCHAR(255),
    fecha_hechos VARCHAR(255),
    descripcion TEXT,
    unidad VARCHAR(255),
    nombre VARCHAR(255),
    primer_apellido VARCHAR(255),
    segundo_apellido VARCHAR(255),
    correo VARCHAR(255),
    boleta VARCHAR(50),
    nombre_tutor VARCHAR(255),
    parentesco VARCHAR(100),
    correo_tutor VARCHAR(255),
    telefono_tutor VARCHAR(20)
);


CREATE INDEX idx_folios_codigo ON folios(codigo_folio);


CREATE INDEX idx_folios_boleta ON folios(boleta);
\end{lstlisting}

% =========================================================
% MODULO DE GESTION DE EVIDENCIAS (ARCHIVOS)
% =========================================================

\section{M\'odulo de Gesti\'on de Evidencias}
Este m\'odulo se encarga de la administraci\'on de los archivos digitales adjuntos a cada expediente. El diseño permite la persistencia de metadatos cr\'iticos y la vinculaci\'on f\'isica de los documentos almacenados en el servidor con su registro en la base de datos.

\subsection{Diccionario de Datos: Tabla Evidencias}
La tabla \texttt{evidencias} registra la ubicaci\'on y caracter\'isticas de los archivos, asegurando que cada documento pueda ser recuperado y validado mediante su tipo y nombre original.

\begin{table}[H]
\centering
\small
\begin{tabularx}{\textwidth}{|l|l|l|X|}
\hline
\rowcolor{azulESCOM} 
\color{white}\textbf{Campo} & \color{white}\textbf{Tipo de Dato} & \color{white}\textbf{Restricciones} & \color{white}\textbf{Descripci\'on} \\ \hline
id & BIGINT & PK, IDENTITY & Identificador \'unico de la evidencia. \\ \hline
nombre\_archivo & VARCHAR(255) & NOT NULL & Nombre original del documento cargado. \\ \hline
tipo\_archivo & VARCHAR(100) & - & Extensi\'on o tipo MIME (PDF, PNG, JPG). \\ \hline
ruta\_almacenamiento & VARCHAR(500) & NOT NULL & Trayectoria f\'isica o URI en el sistema de archivos. \\ \hline
codigo\_folio & VARCHAR(50) & - & Referencia textual del folio (redundancia controlada). \\ \hline
folio\_id & BIGINT & FK, NOT NULL & Relaci\'on con la entidad principal \texttt{folios}. \\ \hline
\end{tabularx}
\caption{Diccionario de datos de la entidad Evidencias.}
\label{tab:dict_evidencias}
\end{table}

\subsection{Definici\'on de la Entidad (JPA)}
La clase \texttt{EvidenciaEntity} implementa el mapeo necesario para gestionar archivos de forma din\'amica. Se destaca la separaci\'on entre el nombre del archivo y su ruta de almacenamiento para facilitar migraciones de servidor.

\begin{lstlisting}[language=Java, caption={Entidad Java para la gesti\'on de metadatos de archivos de evidencia.}, label={lst:java_evidencia}]
@Entity
@Table(name = "evidencias")
@Data
@NoArgsConstructor
@AllArgsConstructor
public class EvidenciaEntity {
    @Id
    @GeneratedValue(strategy = GenerationType.IDENTITY)
    private Long id;

    @Column(nullable = false)
    private String nombreArchivo;
    
    private String tipoArchivo; 
    
    @Column(nullable = false)
    private String rutaAlmacenamiento; 
    
    private String codigoFolio;

    @ManyToOne
    @JoinColumn(name = "folio_id", nullable = false)
    private FolioEntity folio;
}
\end{lstlisting}

\subsection{Definici\'on del Esquema SQL}
El script DDL define una ruta de almacenamiento amplia para prevenir errores por trayectorias de archivos extensas y establece la integridad referencial con la tabla de folios.

\begin{lstlisting}[language=SQL, caption={Script de creaci\'on para la tabla de evidencias digitales.}, label={lst:sql_evidencias}]
CREATE TABLE evidencias (
    id BIGSERIAL PRIMARY KEY,
    nombre_archivo VARCHAR(255) NOT NULL,
    tipo_archivo VARCHAR(100),
    ruta_almacenamiento VARCHAR(500) NOT NULL,
    codigo_folio VARCHAR(50),
    folio_id BIGINT NOT NULL,
    CONSTRAINT fk_folio_evidencia 
        FOREIGN KEY (folio_id) 
        REFERENCES folios(id) 
        ON DELETE CASCADE
);


CREATE INDEX idx_evidencias_folio_id ON evidencias(folio_id);


CREATE INDEX idx_evidencias_codigo_folio ON evidencias(codigo_folio);
\end{lstlisting}

% =========================================================
% MODULO DE QUEJAS (CONSULTA Y SEGUIMIENTO)
% =========================================================

\section{M\'odulo de Quejas y Seguimiento}
Este m\'odulo constituye el n\'ucleo de trazabilidad del sistema, permitiendo el almacenamiento de las quejas formales y la gesti\'on de su ciclo de vida a trav\'es de un historial cronol\'ogico de hitos.

\subsection{Diccionario de Datos: Tabla Quejas}
La tabla \texttt{quejas} utiliza el folio administrativo como identificador \'unico, lo que facilita la interoperabilidad con los procesos f\'isicos de la Defensor\'ia.

\begin{table}[H]
\centering
\small
\begin{tabularx}{\textwidth}{|l|l|l|X|}
\hline
\rowcolor{azulESCOM} 
\color{white}\textbf{Campo} & \color{white}\textbf{Tipo de Dato} & \color{white}\textbf{Restricciones} & \color{white}\textbf{Descripci\'on} \\ \hline
folio & VARCHAR(255) & PK & Identificador alfanum\'erico \'unico del expediente. \\ \hline
fecha\_inicio & TIMESTAMP & DEFAULT NOW() & Fecha y hora de registro inicial de la queja. \\ \hline
asunto & VARCHAR(255) & NOT NULL & Resumen o categor\'ia del motivo de la queja. \\ \hline
status & VARCHAR(50) & DEFAULT 'PENDIENTE' & Estado actual del proceso (PENDIENTE, EN\_PROCESO, RESUELTO). \\ \hline
evidencias & TEXT & - & Referencias o rutas a los archivos probatorios adjuntos. \\ \hline
\end{tabularx}
\caption{Diccionario de datos de la entidad Quejas.}
\label{tab:dict_quejas}
\end{table}

\subsection{Definici\'on de la Entidad (JPA)}
La implementaci\'on en Java destaca por el uso de \texttt{orphanRemoval = true}, garantizando que los hitos del historial se gestionen \'integralmente a trav\'es de la entidad ra\'iz de la queja.

\begin{lstlisting}[language=Java, caption={Entidad Java para la gesti\'on de quejas y expedientes.}, label={lst:java_quejas}]
@Entity
@Table(name = "quejas")
@Data
@NoArgsConstructor
@AllArgsConstructor
public class QuejaEntity {
    @Id
    private String folio;
    
    private LocalDateTime fechaInicio;
    
    @Column(nullable = false)
    private String asunto;
    
    private String status;
    
    @Column(columnDefinition = "TEXT")
    private String evidencias; 
    
    @OneToMany(mappedBy = "queja", cascade = CascadeType.ALL, orphanRemoval = true)
    @OrderBy("fechaHito ASC") 
    private List<HitoEntity> historial = new ArrayList<>();
    
    
    public void addHito(HitoEntity hito) {
        historial.add(hito);
        hito.setQueja(this);
    }
    
    public void removeHito(HitoEntity hito) {
        historial.remove(hito);
        hito.setQueja(null);
    }
}
\end{lstlisting}

\subsection{Definici\'on del Esquema SQL}
El script refleja la capacidad de almacenamiento extendido para el campo de evidencias y la estructura b\'asica de la tabla principal de quejas.

\begin{lstlisting}[language=SQL, caption={Script de creaci\'on para la tabla de quejas.}, label={lst:sql_quejas}]
CREATE TABLE quejas (
    folio VARCHAR(255) PRIMARY KEY,
    fecha_inicio TIMESTAMP DEFAULT CURRENT_TIMESTAMP,
    asunto VARCHAR(255) NOT NULL,
    status VARCHAR(50) DEFAULT 'PENDIENTE',
    evidencias TEXT
);


CREATE INDEX idx_quejas_status ON quejas(status);


CREATE INDEX idx_quejas_fecha_inicio ON quejas(fecha_inicio);
\end{lstlisting}

% =========================================================
% MODULO DE HITOS (HISTORIAL DE SEGUIMIENTO)
% =========================================================

\section{M\'odulo de Hitos de Seguimiento}
Para garantizar la trazabilidad de cada expediente, el sistema implementa una entidad de hitos. Esta permite registrar cronol\'ogicamente cada avance, cambio de estado o acci\'on administrativa realizada sobre una queja espec\'ifica.

\subsection{Diccionario de Datos: Tabla Queja Hitos}
La tabla \texttt{queja\_hitos} mantiene una relaci\'on de muchos a uno con la tabla de quejas, permitiendo que un solo expediente cuente con m\'ultiples registros de seguimiento.

\begin{table}[H]
\centering
\small
\begin{tabularx}{\textwidth}{|l|l|l|X|}
\hline
\rowcolor{azulESCOM} 
\color{white}\textbf{Campo} & \color{white}\textbf{Tipo de Dato} & \color{white}\textbf{Restricciones} & \color{white}\textbf{Descripci\'on} \\ \hline
id & BIGINT & PK, IDENTITY & Identificador \'unico del hito. \\ \hline
descripcion & VARCHAR(255) & NOT NULL & Detalle del evento o acci\'on realizada. \\ \hline
fecha\_hito & TIMESTAMP & NOT NULL & Fecha y hora en que ocurri\'o el evento. \\ \hline
folio\_id & VARCHAR(255) & FK, NOT NULL & Referencia al folio de la queja asociada. \\ \hline
\end{tabularx}
\caption{Diccionario de datos de la entidad Hitos (Seguimiento).}
\label{tab:dict_hitos}
\end{table}

\subsection{Definici\'on de la Entidad (JPA)}
La clase \texttt{HitoEntity} utiliza la anotaci\'on \texttt{@ManyToOne} para establecer la relaci\'on inversa con la entidad de queja. El uso de \texttt{LocalDateTime} asegura precisi\'on en la l\'inea de tiempo del expediente.

\begin{lstlisting}[language=Java, caption={Entidad Java para el registro de hitos de seguimiento.}, label={lst:java_hito}]
@Entity
@Table(name = "queja_hitos")
@Data
@NoArgsConstructor
@AllArgsConstructor
public class HitoEntity {
    @Id
    @GeneratedValue(strategy = GenerationType.IDENTITY)
    private Long id;

    @Column(nullable = false)
    private String descripcion; 
    
    @Column(nullable = false)
    private LocalDateTime fechaHito;

    @ManyToOne
    @JoinColumn(name = "folio_id", nullable = false)
    private QuejaEntity queja;
}
\end{lstlisting}

\subsection{Definici\'on del Esquema SQL}
A continuaci\'on se detalla el script DDL para la creaci\'on de la tabla de hitos, destacando la vinculaci\'on mediante el campo \texttt{folio\_id}.

\begin{lstlisting}[language=SQL, caption={Script de creaci\'on para la tabla de hitos de seguimiento.}, label={lst:sql_hitos}]
CREATE TABLE queja_hitos (
    id BIGSERIAL PRIMARY KEY,
    descripcion VARCHAR(255) NOT NULL,
    fecha_hito TIMESTAMP NOT NULL,
    folio_id VARCHAR(255) NOT NULL,
    CONSTRAINT fk_queja_hito 
        FOREIGN KEY (folio_id) 
        REFERENCES quejas(folio) 
        ON DELETE CASCADE
);


CREATE INDEX idx_hitos_folio_id ON queja_hitos(folio_id);


CREATE INDEX idx_hitos_fecha ON queja_hitos(fecha_hito);
\end{lstlisting}

% =========================================================
% MODULO DE CATALOGOS DEL SISTEMA
% =========================================================

\section{M\'odulo de Cat\'alogos del Sistema}
Para estandarizar la clasificaci\'on de la informaci\'on y mantener la integridad de los datos, se implementan tablas de cat\'alogo con valores predefinidos.

\subsection{Cat\'alogo de Estatus de Queja}

\begin{table}[H]
\centering
\small
\begin{tabularx}{\textwidth}{|l|l|l|X|}
\hline
\rowcolor{azulESCOM} 
\color{white}\textbf{Campo} & \color{white}\textbf{Tipo de Dato} & \color{white}\textbf{Restricciones} & \color{white}\textbf{Descripci\'on} \\ \hline
id & SERIAL & PK & Identificador \'unico del estatus. \\ \hline
nombre & VARCHAR(50) & UNIQUE, NOT NULL & Nombre corto del estado. \\ \hline
descripcion & VARCHAR(255) & - & Explicaci\'on detallada del estado. \\ \hline
\end{tabularx}
\caption{Cat\'alogo de estatus de queja.}
\label{tab:dict_cat_estatus}
\end{table}

\begin{lstlisting}[language=Java, caption={Entidad Java para cat\'alogo de estatus.}, label={lst:java_cat_estatus}]
@Entity
@Table(name = "cat_estatus_queja")
@Data
@NoArgsConstructor
public class CatEstatusQueja {
    @Id
    @GeneratedValue(strategy = GenerationType.IDENTITY)
    private Integer id;
    
    @Column(unique = true, nullable = false, length = 50)
    private String nombre;
    
    @Column(length = 255)
    private String descripcion;
}
\end{lstlisting}

\begin{lstlisting}[language=SQL, caption={Script de creaci\'on e inicializaci\'on del cat\'alogo de estatus.}, label={lst:sql_cat_estatus}]
CREATE TABLE cat_estatus_queja (
    id SERIAL PRIMARY KEY,
    nombre VARCHAR(50) UNIQUE NOT NULL,
    descripcion VARCHAR(255)
);

INSERT INTO cat_estatus_queja (nombre, descripcion) VALUES 
('PENDIENTE', 'Queja registrada pendiente de revisi\'on inicial'),
('EN_PROCESO', 'Queja aceptada y asignada para an\'alisis y seguimiento'),
('RESUELTO', 'Queja con resoluci\'on final o cierre administrativo'),
('RECHAZADO', 'Queja no procedente seg\'un criterios establecidos');
\end{lstlisting}

\subsection{Cat\'alogo de Temas o Categor\'ias}

\begin{table}[H]
\centering
\small
\begin{tabularx}{\textwidth}{|l|l|l|X|}
\hline
\rowcolor{azulESCOM} 
\color{white}\textbf{Campo} & \color{white}\textbf{Tipo de Dato} & \color{white}\textbf{Restricciones} & \color{white}\textbf{Descripci\'on} \\ \hline
id & SERIAL & PK & Identificador \'unico del tema. \\ \hline
nombre & VARCHAR(100) & UNIQUE, NOT NULL & Nombre de la categor\'ia tem\'atica. \\ \hline
descripcion & VARCHAR(255) & - & Descripci\'on detallada del tema. \\ \hline
\end{tabularx}
\caption{Cat\'alogo de temas.}
\label{tab:dict_cat_tema}
\end{table}

\begin{lstlisting}[language=Java, caption={Entidad Java para cat\'alogo de temas.}, label={lst:java_cat_tema}]
@Entity
@Table(name = "cat_tema")
@Data
@NoArgsConstructor
public class CatTema {
    @Id
    @GeneratedValue(strategy = GenerationType.IDENTITY)
    private Integer id;
    
    @Column(unique = true, nullable = false, length = 100)
    private String nombre;
    
    @Column(length = 255)
    private String descripcion;
}
\end{lstlisting}

\begin{lstlisting}[language=SQL, caption={Script de creaci\'on e inicializaci\'on del cat\'alogo de temas.}, label={lst:sql_cat_tema}]
CREATE TABLE cat_tema (
    id SERIAL PRIMARY KEY,
    nombre VARCHAR(100) UNIQUE NOT NULL,
    descripcion VARCHAR(255)
);

INSERT INTO cat_tema (nombre, descripcion) VALUES 
('Acoso escolar', 'Conductas de hostigamiento entre pares dentro del \'ambito escolar'),
('Discriminaci\'on', 'Trato diferenciado o excluyente por motivos de g\'enero, origen, etc.'),
('Abuso de autoridad', 'Ejercicio irregular del poder por parte de personal institucional'),
('Irregularidades acad\'emicas', 'Fallas en procesos de evaluaci\'on, calificaci\'on o administrativos'),
('Violencia de g\'enero', 'Actos de violencia basados en la identidad o expresi\'on de g\'enero'),
('Incumplimiento de funciones', 'Omisi\'on o falta de respuesta por parte del personal');
\end{lstlisting}

\subsection{Cat\'alogo de Tipos de Documento}

\begin{table}[H]
\centering
\small
\begin{tabularx}{\textwidth}{|l|l|l|X|}
\hline
\rowcolor{azulESCOM} 
\color{white}\textbf{Campo} & \color{white}\textbf{Tipo de Dato} & \color{white}\textbf{Restricciones} & \color{white}\textbf{Descripci\'on} \\ \hline
id & SERIAL & PK & Identificador \'unico del tipo de documento. \\ \hline
nombre & VARCHAR(50) & UNIQUE, NOT NULL & Nombre del tipo de documento. \\ \hline
descripcion & VARCHAR(255) & - & Descripci\'on del prop\'osito del documento. \\ \hline
\end{tabularx}
\caption{Cat\'alogo de tipos de documento.}
\label{tab:dict_cat_tipo_doc}
\end{table}

\begin{lstlisting}[language=Java, caption={Entidad Java para cat\'alogo de tipos de documento.}, label={lst:java_cat_tipo_doc}]
@Entity
@Table(name = "cat_tipo_documento")
@Data
@NoArgsConstructor
public class CatTipoDocumento {
    @Id
    @GeneratedValue(strategy = GenerationType.IDENTITY)
    private Integer id;
    
    @Column(unique = true, nullable = false, length = 50)
    private String nombre;
    
    @Column(length = 255)
    private String descripcion;
}
\end{lstlisting}

\begin{lstlisting}[language=SQL, caption={Script de creaci\'on e inicializaci\'on del cat\'alogo de tipos de documento.}, label={lst:sql_cat_tipo_doc}]
CREATE TABLE cat_tipo_documento (
    id SERIAL PRIMARY KEY,
    nombre VARCHAR(50) UNIQUE NOT NULL,
    descripcion VARCHAR(255)
);

INSERT INTO cat_tipo_documento (nombre, descripcion) VALUES 
('Queja formal', 'Documento que expresa inconformidad por actos u omisiones'),
('Denuncia', 'Se\~nalamiento de faltas o irregularidades graves'),
('Solicitud de orientaci\'on', 'Petici\'on de asesor\'ia o informaci\'on sobre derechos'),
('Escrito libre', 'Comunicaci\'on sin formato predeterminado'),
('Comparecencia', 'Declaraci\'on realizada de manera presencial');
\end{lstlisting}

% =========================================================
% NOTA SOBRE ORDEN DE EJECUCION SQL
% =========================================================

\subsection{Orden de Ejecuci\'on de Scripts}
Para garantizar la integridad referencial y evitar errores de dependencias circulares, los scripts DDL deben ejecutarse en el siguiente orden:

\begin{enumerate}
    \item \texttt{cat\_estatus\_queja} (Cat\'alogos)
    \item \texttt{cat\_tema} (Cat\'alogos)
    \item \texttt{cat\_tipo\_documento} (Cat\'alogos)
    \item \texttt{usuarios} (M\'odulo base)
    \item \texttt{folios} (Entidad maestra)
    \item \texttt{quejas} (M\'odulo de quejas)
    \item \texttt{evidencias} (M\'odulo de evidencias - depende de \texttt{folios})
    \item \texttt{queja\_hitos} (M\'odulo de hitos - depende de \texttt{quejas})
    \item \texttt{notificaciones} (M\'odulo de notificaciones - depende de \texttt{usuarios})
    \item \texttt{usuario\_folios} (Tabla auxiliar - depende de \texttt{usuarios})
\end{enumerate}

\begin{center}
\fbox{\parbox{0.9\textwidth}{\centering\textbf{Nota:} Todas las tablas incluyen \'indices adicionales para optimizar el rendimiento de consultas frecuentes. Se recomienda evaluar peri\'odicamente el uso de estos \'indices y ajustarlos seg\'un los patrones de acceso reales del sistema.}}
\end{center}

% =========================================================
% FIN DEL DOCUMENTO
% =========================================================